\documentclass{bmstu}
\usepackage{pdfpages}
\bibliography{parts/biblio}

\setlist[enumerate]{label=\arabic*)}
\setlist[itemize]{label=---}

\usepackage{ragged2e}
\usepackage{array}
\usepackage{tabularx}
\usepackage{adjustbox}  % Добавить в преамбулу

\usepackage{makecell}
\renewcommand\theadalign{c} % Выравнивание по центру и низу
\renewcommand\cellalign{bc}  % Выравнивание содержимого ячеек

\begin{document}
	
	\makethesistitle
		{Информатика и системы управления} % Название факультета
		{Программное обеспечение ЭВМ и информационные технологии} % Название кафедры
		{Сравнительный анализ методов оптимизации архитектуры IFRNet для интерполяции видеокадров с целью повышения вычислительной эффективности} % Тема работы
		{} % Номер группы/ФИО студента (если авторов несколько, их необходимо разделить запятой)
		{} % ФИО научного руководителя
		{} % ФИО консультанта (необязательный аргумент; если консультантов несколько, их необходимо разделить запятой)
		{ }
	
	% Новый титул из ворда
%	\includepdf[pages=-]{title}
	
	\setcounter{page}{2}
%	\include{parts/01-essay}
	\maketableofcontents
	%\cftsetindents{chapter}{\parindent}{12.5mm}
	%\cftsetindents{section}{\parindent}{12.5mm}
	%\cftsetindents{subsection}{\parindent}{12.5mm}
	%\setlength{\cftbeforetoctitleskip}{-22pt}
	%\setlength{\cftaftertoctitleskip}{0pt}
	%\renewcommand{\cftaftertoctitle}{\hfill}
	%\renewcommand{\contentsname}{\hfill\LARGE СОДЕРЖАНИЕ\hfill}
	%\tableofcontents
%	\include{parts/02-def}
	%\include{parts/03-abbr}
	\chapter*{ВВЕДЕНИЕ}
\addcontentsline{toc}{chapter}{ВВЕДЕНИЕ}

Интерполяция видеокадров (Video Frame Interpolation, VFI)~---~это синтез промежуточных кадров между двумя соседними кадрами видеопоследовательности. Технология применяется для повышения частоты кадров, формирования эффекта замедленного воспроизведения и сокращения объёма данных при сжатии видео.

Нейросетевые методы интерполяции обеспечивают высокое качество синтеза, однако точное восстановление промежуточного кадра остаётся затруднительным при нелинейном движении объектов и наличии окклюзий. Современные архитектуры, в частности IFRNet, достигают высоких показателей качества, но характеризуются значительной вычислительной сложностью. Обработка изображений в цветовом пространстве RGB сопряжена с высокими затратами ресурсов, что ограничивает применение таких моделей в системах реального времени.

В связи с этим разработка методов оптимизации, снижающих вычислительную нагрузку без существенной потери визуального качества, в том числе с учётом особенностей зрительного восприятия человека, представляет собой актуальную задачу.

Целью работы является сравнительный анализ методов оптимизации архитектуры IFRNet, направленных на повышение вычислительной эффективности.

Для достижения поставленной цели необходимо решить следующие задачи:
\begin{enumerate}
	\item провести анализ архитектуры IFRNet и обосновать её выбор в качестве объекта оптимизации среди современных методов интерполяции видеокадров;
	\item реализовать методы оптимизации: равномерное сокращение числа каналов (RGB/2) и асимметричную обработку в цветовом пространстве YUV420;
	\item обучить модели на наборе данных Vimeo90K с использованием функций потерь L1, Ternary и Distill;
	\item провести сравнительный анализ методов по показателям вычислительной эффективности (FLOPs, число параметров, FPS) и качества синтеза (PSNR, SSIM, LPIPS, DISTS).
\end{enumerate}

	\chapter{Обзор литературы}

\section{Методы интерполяции видеокадров}

Интерполяция видеокадров (Video Frame Interpolation, VFI) представляет собой задачу синтеза промежуточных кадров между двумя исходными для повышения частоты кадров и плавности видеопоследовательности~\cite{kye2026acevfi}. Данная технология находит широкое применение в различных областях: от создания эффекта замедленного движения и повышения качества видеоконтента до сжатия видео и синтеза новых ракурсов~\cite{huang2022}.

С развитием глубокого обучения методы VFI эволюционировали от традиционных подходов, основанных на компенсации движения, к сложным нейросетевым архитектурам~\cite{kye2026acevfi}. 
Современные методы интерполяции можно классифицировать на несколько категорий в зависимости от архитектурных принципов~\cite{kye2026acevfi}.

\subsection*{Методы на основе оптического потока (Flow-based Methods)}

Методы данной категории основываются на явной оценке оптического потока~---~плотного поля векторов, описывающего смещение каждого пикселя между кадрами~\cite{IFRNet2022}.

Графически оптический поток представляется в виде векторного поля, где каждому пикселю сопоставлен вектор смещения. 
Пример такой визуализации показан на~рис.~\ref{img:flow}.
\includeimage
	{flow}
	{f}
	{H}
	{\textwidth}
	{Визуализация оптического потока $\mathbf{V}_{0\to1}$ как векторного поля.
		Стрелки на рис. (c), наложенные на кадр $I_0$ из (a), показывают смещение каждого пикселя к его положению в кадре $I_1$ (рис. (b)). 
		Длина стрелки пропорциональна величине смещения, направление~---~траектории движения.}
	
Типичный конвейер таких методов включает три основных этапа:

\begin{enumerate}
	\item \textbf{оценка потока}: расчет оптического потока между входными кадрами $I_0$ и $I_1$ с использованием нейросетевых архитектур, таких как FlowNet~\cite{ilg2017}, PWC-Net~\cite{sun2018} или FastFlowNet. Современные подходы часто применяют пирамидальные представления признаков для обработки больших перемещений~\cite{sun2018};
	\item \textbf{деформация (warping)}: использование рассчитанных полей потока для пространственного выравнивания пикселей или признаков входных кадров относительно позиции целевого промежуточного кадра, как правило, с использованием обратного деформирования (backward warping)~\cite{jiang2018};
	\item \textbf{синтез и уточнение}: поскольку деформированные кадры содержат артефакты (разрывы, окклюзии), применяется отдельная сеть синтеза для восстановления целевого кадра~\cite{kye2026acevfi}.
\end{enumerate}

Ключевым преимуществом потоковых методов является явное моделирование движения, обеспечивающее интерпретируемость и возможность интерполяции в произвольные моменты времени~\cite{huang2022}. 
Однако качество таких подходов критически зависит от точности оценки потока и значительно страдает при окклюзиях, больших перемещениях и изменении освещения~\cite{IFRNet2022}.

\subsection*{Методы на основе адаптивных ядер (Kernel-based Methods)}

В отличие от потоковых подходов, методы на основе адаптивных ядер не оценивают движение явно~\cite{niklaus2017}.
Вместо этого они используют архитектуру энкодер-декодер для извлечения признаков из входных кадров и генерации весов для локальных сверточных ядер, которые непосредственно применяются для деформации входных кадров~\cite{niklaus2017}. 
Классическими представителями данного подхода являются SepConvNet~\cite{niklaus2017}, AdaCoFNet и EDSC.
Основное преимущество таких методов~---~устойчивость в слаботекстурных областях, однако они сложны в настройке для произвольной временной интерполяции и диагностики ошибок~\cite{kye2026acevfi}.

\subsection*{Гибридные методы (Hybrid Methods)}

Гибридные подходы комбинируют преимущества потоковых и kernel-методов, используя оценку потока для глобального выравнивания и адаптивные ядра для уточнения локальных деталей~\cite{kye2026acevfi}. 
Это позволяет достигать более высокого качества, но зачастую сопровождается увеличением вычислительной сложности~\cite{huang2022}. 
Примером такого подхода является модель DAIN, использующая глубину для улучшения интерполяции.

\subsection*{Методы на основе трансформеров (Transformer-based Methods)}

В последние годы значительное внимание привлекают методы на основе трансформеров, которые используют механизмы глобального внимания для захвата долгосрочных пространственно-временных зависимостей между пикселями кадров~\cite{kye2026acevfi}. 
В отличие от сверточных сетей, трансформеры способны эффективно моделировать большие перемещения, поскольку операция внимания не ограничена локальным рецептивным полем. 
Среди таких подходов можно выделить VideoINR, использующий трансформеры для неявного представления видео, и TTVFI, комбинирующий трансформеры с модулями уточнения движения. 
Несмотря на высокое качество синтеза, вычислительная сложность трансформеров, квадратично зависящая от размера входного патча, остается существенным ограничением для их применения в реальном времени~\cite{kye2026acevfi}.

\section{Архитектура IFRNet}

Особый интерес для данной работы представляет архитектура \texttt{IFRNet} (Intermediate Feature Refine Network)~\cite{IFRNet2022}. 
Её ключевая инновация заключается в объединении этапов оценки потока и уточнения признаков в единую структуру энкодер-декодер, что позволяет достичь взаимного улучшения этих двух элементов и делает модель лёгкой и быстрой~\cite{IFRNet2022}.

\subsection*{Ключевая идея}

Существующие потоковые методы VFI, достигающие высокого качества, как правило, следуют конвейеру, разделяющему оценку промежуточного потока и синтез изображения на отдельные модули~\cite{IFRNet2022}. 
Это приводит к двум основным недостаткам~\cite{IFRNet2022}:
\begin{enumerate}
	\item разделение промежуточного потока и контекстных признаков на отдельные энкодер-декодеры игнорирует возможность их взаимного улучшения;
	\item каскадная архитектура существенно увеличивает задержку вывода и количество параметров модели, ограничивая применение в реальном времени и на мобильных устройствах~\cite{IFRNet2022}.
\end{enumerate}

\texttt{IFRNet} впервые объединяет эти два этапа в единую модель энкодер-декодер, где промежуточный поток и промежуточные признаки совместно уточняются, помогая друг другу~\cite{IFRNet2022}.
На рис.~\ref{img:diff_arch} представлены различные подходы VFI, основанные на потоке.
\includeimage
	{diff_arch}
	{f}
	{H}
	{\textwidth}
	{Различные подходы VFI, основанные на потоке~\cite{IFRNet2022}}

\subsection*{Структура сети}

\subsubsection*{Пирамидальный энкодер}

Сеть принимает на вход два кадра $I_0$ и $I_1$. 
Общий (shared) энкодер извлекает пирамиду признаков из каждого кадра. 
Энкодер построен из блоков двух сверток $3\times3$ на каждом уровне пирамиды (с шагами 2 и 1 соответственно). 
\texttt{IFRNet} извлекает 4 уровня пирамиды признаков, что составляет 8 сверточных слоев, каждый из которых сопровождается активацией PReLU. 
Количество каналов увеличивается с уровнем: 32, 48, 72 и 96, формируя признаки $\phi_0^k, \phi_1^k$ для каждого уровня $k \in \{1,2,3,4\}$~\cite{IFRNet2022}.

\subsubsection*{Декодеры <<от грубого к точному>>}

Декодеры последовательно обрабатывают признаки от самого грубого (низкое разрешение) уровня к самому детальному. 
Каждый декодер $\mathcal{D}^k$ одновременно предсказывает два компонента~\cite{IFRNet2022}:

\begin{itemize}
	\item \textbf{двусторонний промежуточный поток} ($F_{t \to 0}^k, F_{t \to 1}^k$): векторы движения от целевого момента времени к первому и второму входным кадрам;
	\item \textbf{реконструированный промежуточный признак} ($\hat{\phi}_t^k$): представление целевого кадра в пространстве признаков.
\end{itemize}

Схема архитектуры и декодера представлены на рисунках~\ref{img:ifrnet_arch.png}~-~\ref{img:ifrnet_arch.pdf}.

\includeimage
{ifrnet_arch.png}
{f}
{H}
{.8\textwidth}
{Обзор архитектуры и функций потери данных в \texttt{IFRNet}~\cite{IFRNet2022}}

\includeimage
{ifrnet_arch.pdf}
{f}
{H}
{.8\textwidth}
{Подробная схема декодера на каждом уровне пирамиды~\cite{IFRNet2022}}

Процесс устроен циклически: более точный поток позволяет лучше деформировать исходные пирамидальные признаки посредством обратного деформирования (backward warping), порождая более качественные $\tilde{\phi}_0^k, \tilde{\phi}_1^k$, что дает лучший материал для реконструкции $\hat{\phi}_t^k$. 
В свою очередь, наличие качественного промежуточного признака предоставляет недостающую опорную информацию для более точного предсказания движения на следующем уровне~\cite{IFRNet2022}.

Каждый декодер представляет собой компактную сеть из шести сверток $3\times3$ и одной деконволюции $4\times4$ (с шагом $1/2$). 
Для сохранения эффективности при достаточно большом рецептивном поле третья и пятая свертки обновляют только часть каналов. 
Остаточные соединения способствуют распространению информации и совместному уточнению~\cite{IFRNet2022}.

Благодаря постепенно уточняемому промежуточному признаку, содержащему информацию об окклюзиях и глобальном контексте, \texttt{IFRNet} может генерировать маску слияния и компенсировать детали движения без необходимости в дополнительном модуле синтеза или уточнения, в отличие от многих других методов~\cite{IFRNet2022}.

\subsection*{Синтез кадра}

На финальном этапе выходные карты признаков преобразуются в изображение. 
Сеть генерирует маску слияния $M$ и остаток $R$. Итоговый кадр $\hat{I}_t$ синтезируется по формуле:

\begin{align*}
	\hat{I}_t &= M \odot \tilde{I}_0 + (1 - M) \odot \tilde{I}_1 + R, \\
	\tilde{I}_0 &= w(I_0, F_{t\to0}), \quad \tilde{I}_1 = w(I_1, F_{t\to1}),
\end{align*}

где $\tilde{I}_0$ и $\tilde{I}_1$~---~деформированные входные кадры. Маска $M$ определяет вклад каждого кадра (решая проблему окклюзий), а остаток $R$ добавляет детали текстуры, утраченные при деформации. 
Маска слияния адаптивно усредняет области, видимые в обоих кадрах, а остаток обычно активируется на границах движения и контурах объектов~\cite{IFRNet2022}.

\subsection*{Функции потерь}

Для обучения \texttt{IFRNet} используются специальные функции потерь, позволяющие эффективно обучать модель~\cite{IFRNet2022}:

\begin{itemize}
	\item \textbf{задаче-ориентированный дистилляционный лосс} (task-oriented optical flow distillation loss): этот подход адаптивно регулирует в пространстве робастность дистилляционного знания от более тяжелой предобученной teacher-сети, фокусируясь на изучении полезной информации для синтеза кадра и игнорируя потенциально вредный шум~\cite{IFRNet2022};
	\item \textbf{геометрический консистентный лосс} (feature space geometry consistency loss): этот лосс использует признаки из эталонного промежуточного кадра для регуляризации реконструированных промежуточных признаков, сохраняя правильную геометрическую структуру объектов~\cite{IFRNet2022}.
\end{itemize}

Такой подход позволяет \texttt{IFRNet} достигать качества на уровне SOTA по метрикам \texttt{PSNR}/\texttt{SSIM} при высокой скорости работы и малом числе параметров, так как вся логика упакована в один эффективный блок~\cite{IFRNet2022}.

\section{Методы оптимизации нейросетевых архитектур}

Существует несколько подходов к снижению вычислительной сложности моделей глубокого обучения.

\textbf{Сокращение ширины сети}~---~уменьшение числа каналов в сверточных слоях позволяет значительно сократить число параметров и FLOPs, так как вычислительная стоимость сверточного слоя с ядром $k \times k$ приближённо пропорциональна $H W k^2 C_{in} C_{out}$. 
Одновременное уменьшение входных и выходных каналов в два раза теоретически снижает стоимость примерно в четыре раза.

\textbf{Использование динамических сверток}~---~данный подход позволяет адаптивно применять вычислительные ресурсы только к значимым пикселям изображения. 
В работе CALC-VFS~\cite{giuliani2023calc} предлагается внедрение динамических сверток в архитектуру \texttt{IFRNet}, что позволяет достичь снижения вычислительной сложности на 20-40\% при сохранении схожего качества.

\textbf{Переход к альтернативным цветовым пространствам}~---~использование форматов с субдискретизацией цветоразностных компонент (например, YUV420) позволяет сократить объем входных данных в два раза по сравнению с RGB, что потенциально снижает вычислительную нагрузку на этапе обработки цветовых каналов.


\subsection*{Уменьшение числа признаков в два раза}

Данный метод заключается в замене ширины каналов $C_l$ на

\begin{equation}
	C_l' = \left\lfloor \frac{C_l}{2} \right\rfloor
\end{equation}

для каждого уровня пирамиды. Такая модификация затрагивает энкодер, refine-блоки и часть синтеза. 
Архитектурная логика IFRNet при этом сохраняется: остаются те же уровни пирамиды, тот же порядок уточнения от грубого масштаба к детальному, те же потоки, маски и residual-компоненты.
Архитектурная логика \texttt{IFRNet} при этом сохраняется: остаются те же уровни пирамиды, тот же порядок уточнения от грубого масштаба к детальному, те же потоки, маски и residual-компоненты.

Главное преимущество метода~---~простота внедрения. 
Модель по-прежнему принимает два кадра и предсказывает промежуточный кадр, но использует более узкое внутреннее представление. 
Ожидаемый эффект~---~уменьшение числа параметров, FLOPs и объёма промежуточных активаций.

Недостаток метода связан с уменьшением представительной способности сети. Более узкие признаки могут хуже кодировать мелкие текстуры, сложные движения и области с частичными окклюзиями.

\subsection*{Переход к YUV420}

В цветовом пространстве RGB все три канала имеют одинаковое пространственное разрешение, хотя человеческое зрение более чувствительно к яркости, чем к цветоразностной информации. 
Это свойство широко используется в системах видеокомпрессии: стандарты сжатия, такие как HEVC и H.264, ожидают входные данные в форматах YUV с субдискретизацией цветности.

В YUV-представлении изображение разделяется на яркостную компоненту $Y$ и две цветоразностные компоненты $U$ и $V$. 
Формат YUV420 хранит яркостную компоненту $Y$ в полном разрешении, а цветоразностные компоненты $U$ и $V$~---~с понижением разрешения в два раза по горизонтали и вертикали. 
Для кадра размера $H \times W$ RGB содержит $3HW$ значений, тогда как YUV420:
\begin{equation}
HW + \frac{HW}{4} + \frac{HW}{4} = 1.5HW
\end{equation}

значений, что в два раза меньше данных по сравнению с RGB. Это делает YUV420 перспективным форматом для ускорения интерполяции, поскольку объём входных данных существенно сокращается без потери информации, критичной для зрительного восприятия.

\includeimage
{yuv}
{f}
{H}
{.9\textwidth}
{Фотография и её YUV-компоненты}

\subsubsection*{Стратегии применения YUV420 в IFRNet}

В контексте архитектуры IFRNet возможны две стратегии использования YUV420:

\begin{enumerate}
	\item \textbf{раздельная интерполяция компонент}: интерполяция компонент $Y$, $U$ и $V$ производится раздельно, причём для цветоразностных компонент $U$ и $V$ используется сеть с уменьшенным разрешением или отдельные ветви. 
	Такой подход позволяет сократить вычисления на цветовых каналах, так как они обрабатываются на четверти исходного разрешения.
	
	\item \textbf{интерполяция на основе яркости}: оценка движения выполняется по яркостной компоненте $Y$, а цветоразностные компоненты переносятся с использованием найденного поля движения. 
	Этот подход основан на наблюдении, что движение объектов одинаково для яркости и цвета, а цветовая информация более устойчива к небольшим пространственным ошибкам.
\end{enumerate}

Второй подход является более эффективным, поскольку позволяет избежать дополнительных вычислений на цветовых компонентах.

\subsubsection*{Преимущества и риски}

\textbf{Преимущества} перехода к YUV420:
\begin{itemize}
	\item сокращение объёма входных данных в два раза по сравнению с RGB;
	\item потенциальное снижение вычислительной нагрузки на этапе обработки цветовых каналов;
	\item соответствие стандартам видеосжатия, что упрощает интеграцию в существующие системы;
	\item возможность использования аппаратных ускорителей, оптимизированных для работы с YUV420.
\end{itemize}

\textbf{Основные риски} связаны с возможным появлением цветовых артефактов:
\begin{itemize}
	\item неточное восстановление цветоразностных компонент может приводить к смещению оттенков;
	\item размытие цветных границ из-за пониженного разрешения $U$ и $V$;
	\item ошибки в насыщенных областях и на резких цветовых переходах;
	\item потенциальная потеря мелких цветных деталей.
\end{itemize}

\section{Метрики оценки качества}

Оценка методов интерполяции видеокадров проводится с помощью комбинации метрик, измеряющих различные аспекты качества~\cite{kye2026acevfi}. 
Выбор метрик определяется необходимостью комплексного анализа: пиксельные метрики позволяют оценить численную точность восстановления, тогда как перцептуальные метрики отражают визуальное качество, согласованное с человеческим восприятием~\cite{kye2026acevfi}. 
В данном исследовании используются четыре метрики: \texttt{PSNR}, \texttt{SSIM}, \texttt{LPIPS} и \texttt{DISTS}, что позволяет всесторонне сравнить эффективность предлагаемых методов оптимизации.

\subsection*{Пиксельные метрики}

Пиксельные метрики сравнивают интерполированный кадр $\hat{I}_t$ с эталонным $I_t$ на уровне отдельных пикселей или их локальных статистик.

\subsubsection*{PSNR (Peak Signal-to-Noise Ratio)}

\texttt{PSNR} является наиболее распространённой метрикой оценки качества восстановления изображений. Она вычисляется на основе среднеквадратичной ошибки (Mean Squared Error, MSE) между пикселями эталонного и восстановленного изображений:

\begin{equation}
	\text{MSE} = \frac{1}{HW} \sum_{i=1}^{H} \sum_{j=1}^{W} \left( I_t(i,j) - \hat{I}_t(i,j) \right)^2,
\end{equation}

\begin{equation}
	\text{PSNR} = 10 \cdot \log_{10}\left(\frac{MAX_I^2}{\text{MSE}}\right),
\texttt{PSNR} = 10 \cdot \log_{10}\left(\frac{MAX_I^2}{\text{MSE}}\right),
\end{equation}

где $H$ и $W$~---~высота и ширина изображения, $MAX_I$~---~максимальное значение пикселя (для 8-битных изображений $MAX_I = 255$). Высокие значения \texttt{PSNR} указывают на высокое численное сходство между изображениями.

Преимуществом \texttt{PSNR} является её простота, физический смысл и возможность легко сравнивать результаты с другими работами благодаря широкой распространённости. Однако главный недостаток заключается в том, что \texttt{PSNR} плохо коррелирует с человеческим восприятием: два изображения могут иметь одинаковый \texttt{PSNR}, но одно будет выглядеть отлично, а другое~---~содержать заметные артефакты, поскольку \texttt{PSNR} не учитывает структурную информацию и пространственную организацию пикселей.

\subsubsection*{SSIM (Structural Similarity Index)}

\texttt{SSIM} была разработана для преодоления недостатков \texttt{PSNR} и оценки сходства изображений с учётом особенностей зрительной системы человека. Метрика вычисляет сходство по трём компонентам:

\begin{equation}
	\text{SSIM}(x, y) = \frac{(2\mu_x\mu_y + C_1)(2\sigma_{xy} + C_2)}{(\mu_x^2 + \mu_y^2 + C_1)(\sigma_x^2 + \sigma_y^2 + C_2)},
\texttt{SSIM}(x, y) = \frac{(2\mu_x\mu_y + C_1)(2\sigma_{xy} + C_2)}{(\mu_x^2 + \mu_y^2 + C_1)(\sigma_x^2 + \sigma_y^2 + C_2)},
\text{SSIM}(x, y) = \frac{(2\mu_x\mu_y + C_1)(2\sigma_{xy} + C_2)}{(\mu_x^2 + \mu_y^2 + C_1)(\sigma_x^2 + \sigma_y^2 + C_2)},
\end{equation}

где $x$ и $y$~---~соответствующие локальные блоки эталонного и восстановленного изображений; $\mu_x$, $\mu_y$~---~средние значения; $\sigma_x$, $\sigma_y$~---~дисперсии; $\sigma_{xy}$~---~ковариация; $C_1$ и $C_2$~---~константы, вводимые для стабилизации деления. Итоговое значение \texttt{SSIM} вычисляется как среднее по всем локальным блокам изображения.

\texttt{SSIM} значительно лучше согласуется с человеческим восприятием, чем \texttt{PSNR}, и стала стандартом де-факто в области обработки изображений. Однако метрика имеет ограничения: она может давать нестабильные результаты на изображениях с низкой контрастностью или яркостью, а также чувствительна к мелким пространственным сдвигам и деформациям.

\subsection*{Перцептуальные метрики}

Перцептуальные метрики оценивают визуальное качество через признаки, извлечённые из глубоких нейросетей, что позволяет приблизить оценку к субъективному восприятию человека~\cite{kye2026acevfi}.

\subsubsection*{LPIPS (Learned Perceptual Image Patch Similarity)}

\texttt{LPIPS} использует предобученную свёрточную нейросеть (как правило, VGG) для извлечения карт признаков из промежуточных слоёв. Метрика вычисляет взвешенное расстояние между признаками эталонного и восстановленного изображений на каждом уровне:

\begin{equation}
	\text{LPIPS}(I_t, \hat{I}_t) = \sum_{l} \sum_{c} \lambda_{l,c} \cdot \| \phi_{l,c}(I_t) - \phi_{l,c}(\hat{I}_t) \|_2,
\end{equation}

где $\phi_{l,c}$~---~карта признаков на слое $l$ для канала $c$, а $\lambda_{l,c}$~---~изученные веса, оптимизированные для согласования с субъективными оценками человека. Низкие значения \texttt{LPIPS} указывают на лучшее перцептуальное качество.

\texttt{LPIPS} демонстрирует более высокую корреляцию с человеческими оценками по сравнению с \texttt{PSNR} и \texttt{SSIM}, что делает её предпочтительной для задач, где визуальное качество важнее численной точности. Однако метрика не лишена недостатков: она может быть обманута специальными атаками, и её работа не всегда интерпретируема. Кроме того, \texttt{LPIPS} чувствительна к глобальным изменениям цвета и контраста.

\subsubsection*{DISTS (Deep Image Structure and Texture Similarity)}

\texttt{DISTS} является более современной метрикой, разработанной для раздельной оценки структурного и текстурного сходства. В отличие от \texttt{LPIPS}, \texttt{DISTS} вычисляет две компоненты:

\begin{enumerate}
	\item \textbf{Структурное сходство} (Structure Similarity)~---~оценивает пространственное расположение признаков, то есть точную структуру объектов;
	\item \textbf{Текстурное сходство} (Texture Similarity)~---~оценивает статистики признаков по пространству, что позволяет быть устойчивым к передискретизации текстур.
\end{enumerate}

Математически метрика может быть представлена как:

\begin{equation}
	\text{DISTS}(I_t, \hat{I}_t) = \sum_{l} \sum_{c} \lambda_{l,c} \cdot \left\| \text{corr}\left(\phi_{l,c}(I_t)\right) - \text{corr}\left(\phi_{l,c}(\hat{I}_t)\right) \right\|_2,
\end{equation}

где $\text{corr}$ обозначает вычисление корреляционной структуры карты признаков, а $\lambda_{l,c}$~---~изученные веса.

Ключевая особенность \texttt{DISTS} заключается в том, что она позволяет учесть вариативность текстур (например, при синтезе травы, воды или песка), где небольшие перестановки текстурных элементов не должны существенно влиять на оценку качества. Исследования показывают, что \texttt{DISTS} демонстрирует наивысшую корреляцию с субъективными оценками человека среди всех современных метрик и устойчива к небольшим геометрическим искажениям.

\subsection*{Обоснование выбора метрик}

Комплексное использование пиксельных и перцептуальных метрик позволяет объективно оценивать как численную точность, так и визуальное качество интерполированных кадров, что критически важно для сравнения методов оптимизации~\cite{kye2026acevfi}.

В контексте данной работы выбор четырёх метрик (\texttt{PSNR}, \texttt{SSIM}, \texttt{LPIPS}, \texttt{DISTS}) обусловлен следующими соображениями:
\t\item \textbf{\texttt{PSNR}} и \textbf{\texttt{SSIM}} позволяют оценить, насколько сильно изменились пиксели и простая структура изображения после применения методов оптимизации (сокращения каналов или перехода в YUV420);

\begin{itemize}
	\item \textbf{\texttt{PSNR}} и \textbf{\texttt{SSIM}} позволяют оценить, насколько сильно изменились пиксели и простая структура изображения после применения методов оптимизации (сокращения каналов или перехода в YUV420);
	\item \textbf{\texttt{LPIPS}} и \textbf{\texttt{DISTS}} показывают, насколько эти изменения заметны для человеческого глаза, что критически важно, поскольку цель оптимизации~---~сохранить визуальное качество при повышении вычислительной эффективности.
\end{itemize}

Например, если после оптимизации значение \texttt{PSNR} снизилось, но \texttt{LPIPS} и \texttt{DISTS} изменились незначительно, это означает, что выигрыш в скорости достигнут без заметной для наблюдателя потери качества.
	\chapter{Сравнение методов}

Сравнение проводится по двум группам метрик: вычислительной эффективности (FLOPs, число параметров, время инференса) и качеству синтеза (PSNR, SSIM, LPIPS).

\section{Экспериментальная установка}

\subsection*{Модели для сравнения}

В рамках исследования были обучены и протестированы следующие модели на основе архитектуры IFRNet:
\begin{itemize}
	\item \texttt{IFRNet (baseline)}~---~оригинальная архитектура, предложенная в работе \cite{IFRNet2022}. 
	Использует полное цветовое пространство RGB и стандартную функцию потерь: комбинацию L1, Ternary и дистилляционного лосса.
	
	\item \texttt{IFRNet\_RGB2}~---~модификация с равномерным сокращением числа каналов в два раза на всех уровнях пирамиды признаков (каналы 16, 24, 36, 48 вместо 32, 48, 72, 96). 
	Функция потерь идентична базовой модели.
	
	\item \texttt{IFRNet\_YUV420}~---~модификация с переходом к цветовому пространству YUV420 и интерполяцией на основе яркостной компоненты $Y$. 
	Цветоразностные компоненты $U$ и $V$ переносятся с использованием оптического потока, вычисленного по яркости. 
	Функция потерь аналогична базовой модели.
	
	\item \texttt{IFRNet\_RGB2+LPIPS}~---~архитектура \texttt{IFRNet\_RGB2} с добавлением перцептуальной функции потерь LPIPS~\cite{zhang2018} в общую функцию потерь.
	
	\item \texttt{IFRNet\_YUV420+LPIPS}~---~архитектура \texttt{IFRNet\_YUV420} с добавлением перцептуальной функции потерь LPIPS.
\end{itemize}

\subsection*{Метрики оценки}

Оценка производилась по двум категориям метрик:

\textbf{Метрики эффективности:}
\begin{itemize}
	\item \textbf{FLOPs} (Floating Point Operations)~---~количество операций с плавающей запятой, необходимое для обработки одного кадра. Рассчитывалось с помощью библиотеки \texttt{thop} для изображений размером $1920 \times 1088$.
	\item \textbf{Params}~---~количество обучаемых параметров модели (в миллионах).
	\item \textbf{Inference Time}~---~среднее время инференса на один кадр (в секундах).
\end{itemize}

\textbf{Метрики качества:}
\begin{itemize}
	\item \textbf{PSNR} (Peak Signal-to-Noise Ratio)~---~пиксельная метрика, измеряющая отношение сигнал/шум в децибелах.
	\item \textbf{SSIM} (Structural Similarity Index)~---~метрика структурного сходства, учитывающая яркость, контраст и структуру изображения~\cite{wang2004}.
	\item \textbf{LPIPS} (Learned Perceptual Image Patch Similarity)~---~перцептуальная метрика, использующая признаки из предобученной сети VGG для оценки визуального сходства~\cite{zhang2018}.
\end{itemize}

\subsection*{Датасет и условия обучения}

Обучение всех моделей проводилось на датасете Vimeo90K. Входные изображения имели разрешение $1920 \times 1088$ пикселей. 
Для моделей с LPIPS в функции потерь использовался коэффициент $\lambda = 0.1$ для взвешивания перцептуальной компоненты.

\section{Результаты экспериментов}

\subsection*{Количественное сравнение}

Результаты сравнения всех моделей представлены в таблице~\ref{tab:comparison}.

\begin{table}[htbp]
	\centering
	\caption{Сравнение методов оптимизации IFRNet}
	\label{tab:comparison}
	\adjustbox{max width=\textwidth}{
		\begin{tabular}{l c c c c c c c}
			\toprule
			\textbf{Модель} & \texttt{FLOPs (↓)} & \texttt{Params (↓)} & \texttt{Inference Time (↓)} & \texttt{PSNR (↑)} & \texttt{SSIM (↑)} & \texttt{LPIPS (↓)} & \texttt{DISTS (↓)} \\
			& & & \textbf{(с/кадр)} & & & & \\
			\midrule
			\texttt{IFRNet (baseline)} & 479.441G & 4.960M & 0.034 & 35.698 & 0.9790 & 0.0517 & 0.0421 \\
			\texttt{IFRNet\_RGB2} & 126.191G & 1.00M & 0.015 & 34.897 & 0.9756 & 0.0542 & 0.0453 \\
			\texttt{IFRNet\_YUV420} & 122.682G & 1.00M & 0.017 & 34.446 & 0.9737 & 0.0622 & 0.0527 \\
			\texttt{IFRNet\_RGB2+LPIPS} & 126.191G & 1.00M & 0.015 & 34.659 & 0.9743 & 0.0380 & 0.0389 \\
			\texttt{IFRNet\_YUV420+LPIPS} & 122.682G & 1.00M & 0.017 & 35.121 & 0.9761 & 0.0481 & 0.0405 \\
			\bottomrule
		\end{tabular}
	}
\end{table}

\subsection*{Анализ эффективности}

\textbf{Сокращение вычислительной сложности.}
Все предложенные методы оптимизации демонстрируют значительное снижение вычислительной сложности по сравнению с базовой моделью IFRNet. 
Наиболее существенное ускорение достигается за счёт сокращения ширины сети: модели \texttt{IFRNet\_RGB2} и \texttt{IFRNet\_RGB2+LPIPS} показывают снижение FLOPs в 3.8 раза (с 479.441G до 126.191G) и ускорение инференса более чем в 2 раза (с 0.034 до 0.015 секунды на кадр).

Модель \texttt{IFRNet\_YUV420} также показывает значительное снижение FLOPs до 122.682G, что сопоставимо с RGB2-версиями, однако время инференса несколько выше (0.017 с/кадр), что может быть связано с дополнительными операциями преобразования цветовых пространств и деформации цветовых каналов.

\textbf{Число параметров.} 
Сокращение ширины сети в два раза позволило уменьшить число параметров с 4.96M до 1.00M. 
Модели на основе YUV420 сохраняют такое же количество параметров, что объясняется неизменной архитектурой сети (изменения касаются только входных данных и стратегии синтеза цветовых компонент).

\subsection*{Анализ качества}

\textbf{Пиксельные метрики.} 
Все оптимизированные модели показывают некоторое снижение PSNR и SSIM по сравнению с базовой архитектурой. 
Наименьшую потерю качества демонстрирует модель \texttt{IFRNet\_YUV420+LPIPS}, которая достигает PSNR = 35.121 дБ при снижении FLOPs в 3.9 раза. Это подтверждает гипотезу о том, что цветовая информация менее критична для восприятия, и её обработка с пониженным разрешением позволяет сохранить высокое качество синтеза.

Модели без LPIPS в функции потерь (\texttt{IFRNet\_RGB2} и \texttt{IFRNet\_YUV420}) показывают ожидаемое снижение PSNR на 0.8-1.25 дБ, что является приемлемой ценой за значительное ускорение работы модели.

\textbf{Перцептуальное качество.}
Наиболее интересные результаты получены по метрике LPIPS. 
Добавление перцептуальной функции потерь в процесс обучения позволило существенно улучшить визуальное качество генерируемых кадров. 
Модель \texttt{IFRNet\_RGB2+LPIPS} показывает значение LPIPS = 0.0380, что даже лучше, чем у базовой модели (0.0517), при том же уровне вычислительной эффективности, что и \texttt{IFRNet\_RGB2}.

Модель IFRNet\_YUV420+LPIPS также демонстрирует значительное улучшение LPIPS по сравнению с версией без LPIPS (0.0481 против 0.0622), при этом оставаясь близкой к базовой модели по этому показателю.

\subsection*{Сравнительный анализ баланса эффективности и качества}

На рисунке~\ref{fig:pareto} представлена визуализация компромисса между вычислительной эффективностью и качеством синтеза для всех рассмотренных моделей.

% \begin{figure}[htbp]
	%     \centering
	%     \includegraphics[width=\textwidth]{pareto_chart.png}
	%     \caption{График компромисса между FLOPs и качеством (LPIPS) для различных моделей}
	%     \label{fig:pareto}
	% \end{figure}

Можно выделить три группы моделей по соотношению качество/эффективность:

\begin{enumerate}
	\item \textbf{максимальная эффективность}: IFRNet\_RGB2 и IFRNet\_RGB2+LPIPS обеспечивают наибольшее ускорение (в 3.8 раза по FLOPs и в 2.3 раза по времени инференса). При этом модель с LPIPS значительно превосходит версию без LPIPS по перцептуальному качеству.
	
	\item \textbf{сбалансированное решение}: IFRNet\_YUV420+LPIPS предлагает оптимальный баланс между скоростью и качеством, обеспечивая значительное ускорение (FLOPs ↓ в 3.9 раз) при минимальной потере качества по пиксельным метрикам и улучшении перцептуальных характеристик.
	
	\item \textbf{базовая модель}: IFRNet (baseline) демонстрирует наилучшее пиксельное качество, но ценой значительно более высоких вычислительных затрат.
\end{enumerate}

\section{Визуальное сравнение}

Для качественной оценки результатов было проведено визуальное сравнение интерполированных кадров, полученных различными моделями. На рисунках~\ref{fig:visual_comparison1} и~\ref{fig:visual_comparison2} представлены примеры синтезированных кадров для различных сцен.

% \begin{figure}[htbp]
	%     \centering
	%     \includegraphics[width=\textwidth]{visual_comparison1.png}
	%     \caption{Визуальное сравнение интерполяции для сцены с быстрым движением}
	%     \label{fig:visual_comparison1}
	% \end{figure}

% \begin{figure}[htbp]
	%     \centering
	%     \includegraphics[width=\textwidth]{visual_comparison2.png}
	%     \caption{Визуальное сравнение интерполяции для сцены со сложными текстурами}
	%     \label{fig:visual_comparison2}
	% \end{figure}

\textbf{Анализ визуальных результатов:}
\begin{itemize}
	\item Модели без LPIPS (IFRNet\_RGB2, IFRNet\_YUV420) демонстрируют некоторое сглаживание текстур и потерю мелких деталей, что особенно заметно на сложных природных сценах.
	\item Добавление LPIPS в функцию потерь позволяет восстановить резкость и текстуры: модели с LPIPS дают более детализированные и естественные изображения.
	\item Модель IFRNet\_YUV420+LPIPS показывает результаты, визуально близкие к базовой модели, при значительно меньшей вычислительной нагрузке.
\end{itemize}
	\chapter*{ЗАКЛЮЧЕНИЕ}
\addcontentsline{toc}{chapter}{ЗАКЛЮЧЕНИЕ}

В работе предложены и исследованы методы оптимизации архитектуры IFRNet, направленные на повышение вычислительной эффективности при сохранении качества интерполяции видеокадров.

\subsection*{Основные результаты}

\begin{enumerate}
	\item Предложенные методы (сокращение числа каналов и переход к YUV420) обеспечили снижение вычислительной сложности в 3,8--3,9 раза по FLOPs и уменьшение числа параметров в 4,96 раза (с 4,96~M до 1,00~M); ускорение инференса составило более чем в два раза (с 0,032 до 0,014--0,015~с/кадр).

	\item Наилучший баланс между эффективностью и качеством показала модель \texttt{IFRNet\_RGB2+LPIPS}: при снижении FLOPs в 3,8 раза она обеспечивает наилучшее среди всех моделей перцептуальное качество (LPIPS = 0,0361, DISTS = 0,0222), превосходящее в том числе базовую архитектуру (LPIPS = 0,0503, DISTS = 0,0342), при близких пиксельных метриках (PSNR = 35,129~дБ, SSIM = 0,9774).

	\item Добавление перцептуальной функции потерь LPIPS компенсирует снижение визуального качества, вызванное архитектурной оптимизацией: относительно модели \texttt{IFRNet\_RGB2} без LPIPS улучшение составило около 31~\% по метрике LPIPS и около 39~\% по метрике DISTS.

	\item Переход к цветовому пространству YUV420 обеспечивает минимальные значения FLOPs и времени инференса, однако уступает варианту RGB2 по перцептуальным метрикам; его применение целесообразно при интеграции в конвейеры обработки видео, изначально использующие представление YUV.
\end{enumerate}

\subsection*{Перспективы дальнейших исследований}

Дальнейшие исследования могут быть направлены на:
\begin{itemize}
	\item разработку гибридной стратегии, адаптивно выбирающей между режимами RGB2 и YUV420 в зависимости от характеристик сцены (уровня детализации текстур, интенсивности движения), что позволит динамически балансировать вычислительную эффективность и качество интерполяции;
	\item применение других перцептуальных функций потерь (например, DISTS) в процессе обучения для дальнейшего повышения визуального качества оптимизированных моделей;
	\item адаптацию предложенных моделей для мобильных и встраиваемых устройств с использованием аппаратного ускорения (TensorRT, TFLite) и методов квантования.
\end{itemize}

	
	\makebibliography
	
%	\include{parts/10-appendix}
	
\end{document}